\section{Geometric progressions over products of finite fields}
\label{sec:products}
Asada, Fourakis, Manski, McNew, Miller and Moreland ask in Question~6.5
of their \href{https://doi.org/10.1016/j.ffa.2016.10.002}{2017 paper}
how their density problem changes over other finite coefficient rings.
Their Remark~1.1 specifies polynomial ratios that are nonunits.
A finite product of fields gives a new feature: multiplication by $x$ can be
switched on independently in any subset of the factors. This yields a sharp
criterion for whether a positive upper density can survive as factors are added.

For $R_b=\prod_{i=1}^b\mathbb F_{q_i}$ take
$B_n(R_b)=\{f\in R_b[x]:\text{all coefficients above degree }n\text{ vanish}\}$.
It is the product of the coordinate balls and has size $|R_b|^{n+1}$.
Write $\Delta(R_b)$ for the supremum of
$\limsup_{n\to\infty}|E\cap B_n(R_b)|/|B_n(R_b)|$ over sets $E$
avoiding triples $f,rf,r^2f$ with a nonunit $r\in R_b[x]$.
All constructions below have every component nonzero. The upper bound
uses only progressions with three distinct terms. Thus possible conventions
about repeated terms or zero leave the results unchanged: the proportion of
tuples having a zero component is at most $\sum_i q_i^{-(n+1)}$ at fixed $b$.

\begin{theorem}[Criterion for growing products]\label{thm:products}
For an arbitrary sequence of prime powers $q_1,q_2,\ldots$, the limit
$\lim_{b\to\infty}\Delta(R_b)$ exists. It is positive precisely when
$\sum_i q_i^{-2}$ converges. For every finite $b$,
\[
 \prod_{i=1}^b(1-q_i^{-2})\ \le\ \Delta(R_b)
 \ \le\ \mathbb E\, a_K,\qquad
 K=\sum_{i=1}^b X_i,\qquad
 \Pr(X_i=1)=\frac{3}{q_i^2+q_i+1},
\]
where $X_i$ are independent Bernoulli variables and $a_k$ is the largest
proportion of the ternary word cube $\{0,1,2\}^k$ containing no
combinatorial line, with $a_0=1$. In particular $\Delta(\mathbb F_q^b)\to0$
for fixed $q$. There is an absolute $C>0$ such that for every $b\ge1$
and all $q_i$,
\[
 \Delta(R_b)\ge \exp\!\big(-C\sqrt{\log(2+b)}\big).
\]
\end{theorem}

\subsection{Packing the polynomial ball with word cubes}
A combinatorial line has three words, with a nonempty set of variable
positions reading respectively $0,1,2$ in unison and all other positions
fixed. Fix a field of order $q$. Each nonzero polynomial in its ball
has a unique expression $x^eg$ where $g(0)\ne0$. For each fixed $g$,
partition the exponent interval $0\le e\le n-\deg g$ into disjoint
successive triples starting at $0,3,6,\ldots$, leaving up to two singleton
cells. Put zero in a singleton. For $3j+2\le n$, the number of triples
starting at exponent $3j$ equals $(q-1)q^{n-2-3j}$, the number of choices
of $g$ having nonzero constant term and degree at most $n-2-3j$.
Consequently the fraction of the ball lying in triple cells is
\[
 \theta_{q,n}=3(q-1)\sum_{0\le j\le(n-2)/3}q^{-3j-3}
 \ \longrightarrow\ \frac{3}{q^2+q+1}.
\]

Choose a point uniformly in the product ball, then condition on the cell
of each coordinate. The coordinates in triple cells are independent
uniform letters. Their count $K_n$ has the law of a sum of independent
Bernoulli variables with parameters $\theta_{q_i,n}$. Within a cell
with $k$ active coordinates, each combinatorial line is
$f,rf,r^2f$: take $r_i=x$ on its variable positions and $r_i=1$
elsewhere. The three tuples are distinct, since the variable positions
belong to nonzero chains. Conditional occupancy by any admissible set
is at most $a_k$. Averaging gives an upper bound $\mathbb E a_{K_n}$
even for the finite-ball maximum. At fixed $b$, passage to the limit in
this finite distribution gives the asserted upper bound.

The \href{https://annals.math.princeton.edu/2012/175-3/p06}{density
Hales--Jewett theorem} gives $a_k\to0$. Write
$\mu_b=\sum_{i\le b}3/(q_i^2+q_i+1)$. For $q_i\ge2$ the summands
are between $(12/7)q_i^{-2}$ and $3q_i^{-2}$. Thus $\mu_b\to\infty$
when $\sum_iq_i^{-2}$ diverges. Since $\mathrm{Var}(K)\le\mu_b$,
the probability $\Pr(K<\mu_b/2)$ is at most $4/\mu_b$. The contribution
of larger $K$ is bounded by $\sup_{k\ge\mu_b/2}a_k\to0$. This proves the
vanishing assertion, using only the $2^b-1$ possible nonconstant ratios
whose coordinates are in $\{1,x\}$.

\subsection{Passing from a degree cutoff to a single infinite set}
Let $M_b(n)$ be the largest admissible fraction of $B_n(R_b)$.
Then
\begin{equation}\label{eq:glueproduct}
 \Delta(R_b)=\limsup_{n\to\infty} M_b(n).
\end{equation}
For the nontrivial direction choose $n_j$ so that $M_b(n_j)$ approaches
its upper limit, $n_j-j>2n_{j-1}$, and $n_j-2j>n_{j-1}$.
In a maximizing set at cutoff $n_j$ keep only tuples with all degrees
between $n_j-j$ and $n_j$, taking $n_1>1$. This deletes a proportion
at most $\sum_i q_i^{-j-1}$. Form the union of the retained boxes.
In any progression in this union all $f_i,r_i$ are nonzero, and
the three degrees in each coordinate are nondecreasing and equally spaced.
If the last term alone lies in the latest box, each last degree is at
most $2n_{j-1}$, too small. If the last two lie there, the first degree
in each coordinate is at least $n_j-2j>n_{j-1}$. Within one box avoidance
was imposed already. This proves \eqref{eq:glueproduct}. Slicing a finite
product ball by its last coordinate gives $M_{b+1}(n)\le M_b(n)$:
any forbidden triple in a slice lifts by adjoining ratio $1$. Hence
$\Delta(R_b)$ is nonincreasing and has a limit.

The tuples with every component degree in $\{n-1,n\}$, $n\ge1$,
occupy exactly $\prod_i(1-q_i^{-2})$ of $B_n(R_b)$. A surviving nonunit
ratio would have at least one component of positive degree. In that
coordinate its first and last degrees would differ by at least two.
This proves the lower product bound by \eqref{eq:glueproduct}.
Since $0<q_i^{-2}\le1/4$, the infinite product is positive exactly when
the sum of its deficits converges. Together with the cube bound this
completes the dichotomy.

\begin{corollary}[The regime of densities close to one]\label{cor:nearone}
For any sequence of finite products (with no nesting assumption), set
$S=\sum_i q_i^{-2}$ for each product. Along such a sequence,
$\Delta\to0$ exactly when $S\to\infty$, and $\Delta\to1$ exactly when
$S\to0$. In the latter regime the first-order estimate is uniform:
\[
 1-\Delta=S+O(S^{3/2}).
\]
\end{corollary}
\begin{proof}
The upper estimate from word cubes tends to zero when $S\to\infty$
by the variance argument above. The product lower bound is at least
$\exp(-4S/3)$, since $-\log(1-u)\le4u/3$ for $0\le u\le1/4$.
For the other endpoint, slicing any nonempty-dimensional ternary cube
into parallel lines shows $a_k\le2/3$ for $k\ge1$. With
$p_i=3/(q_i^2+q_i+1)$ this gives
\[
 \frac{1-\prod_i(1-p_i)}{3}\le 1-\Delta\le S.
\]
As $\sum_i p_i\ge12S/7$, the left side stays positive when $S$ stays
bounded below. As $S\to0$ each $q_i^{-1}\le\sqrt S$. Thus
$\sum_i p_i=3S+O(S^{3/2})$.
The first two inclusion--exclusion bounds give
$1-\prod_i(1-p_i)=\sum_i p_i+O((\sum_i p_i)^2)$. This proves the
uniform expansion.
\end{proof}

For a finite illustration, the bounds for two copies of $\mathbb F_2$ are
$9/16\le\Delta\le38/49$; for
$\mathbb F_2\times\mathbb F_3\times\mathbb F_5$ they are
$16/25\le\Delta\le322/403$. These use the values
$a_1=a_2=a_3=2/3$ from
\href{https://arxiv.org/abs/1002.0374}{D.~H.~J. Polymath's finite cube calculation}.
The program \texttt{check\_products.py} verifies the chain counts by direct
enumeration at small cutoffs, checks all $37$ lines in a three-coordinate
cell, and computes the displayed bounds in rational arithmetic.

\subsection{Using the repository's Behrend lemma for a lower bound}
\label{sec:behrend-products}
The source lemma
\href{https://github.com/lsernau/grh-lean/blob/72fead102f663d5c7a407964a13dc0cf9d977207/grh/src/New/Reduction/Data/Existence.lean\#L77-L113}{\texttt{exists\_three\_term\_free}}
works with points on a sphere in a digit box. It supplies $S\subset[0,o)$
of size at least $H^d/(dH^2)$ for $(LH)^d\le o$, avoiding nontrivial
weighted three-term congruences with positive weights of sum at most $L$.
In its application to infinite places, multiplicative relations impose
additive constraints on exponents. Here the transfer uses ordinary
midpoints ($L=2$) and projects each polynomial tuple to a single integer.

In particular there is an absolute $C_0$ such that for each $m\ge2$ some
$S\subset\{0,\ldots,m-1\}$ contains no nontrivial three-term arithmetic
progression and has size at least $m\exp(-C_0\sqrt{\log(2+m)})$.
For large $m$ take $d=\lfloor\sqrt{\log_2m}\rfloor$,
$H=\lfloor m^{1/d}/2\rfloor$, and $o=m$. Then $(2H)^d\le m$,
$H^d\ge m/4^d$, and $H^2\le m^{2/d}$. Substitution gives the bound;
a larger absolute $C_0$ covers the finitely many smaller values.

For a random tuple at cutoff $n$ with nonzero components consider its
total degree deficit $D_n=\sum_{i=1}^b(n-\deg f_i)$. For fixed $b$ its
limiting law is the sum $D$ of independent geometric variables with
\[
 \Pr(D_i=t)=(1-q_i^{-1})q_i^{-t}\quad(t\ge0),\qquad
 V=\mathrm{Var}(D)=\sum_{i=1}^b\frac{q_i}{(q_i-1)^2}\le2b .
\]
Indeed the stated geometric probability for a given $t\le n$ is the
exact unconditioned probability of degree $n-t$; zero components have
vanishing probability. An integer interval $I$ of length (cardinality)
$m\le4\sqrt V+3$, enlarged to at least two points if necessary, captures
at least $3/4$ of $D$ by Chebyshev's inequality. As $h$ ranges over
$I-\{0,\ldots,m-1\}$, an interval of $2m-1$ integers, each $t\in I$
belongs to exactly $|S|$ sets $h+S$. Consequently for some $h$,
\[
 \Pr(D\in h+S)\ge\frac{3|S|}{4(2m-1)}.
\]
At cutoff $n$ take the nonzero-component tuples with $D_n\in h+S$.
Along a candidate progression the total degree rises by $\sum_i\deg r_i>0$,
so its three deficits would form a forbidden arithmetic progression.
This admissible finite construction has limiting proportion at least the
last displayed bound. Applying \eqref{eq:glueproduct} and $V\le2b$ proves
the stated lower estimate, after enlarging the absolute constant. In fact
the proof gives $\exp(-C\sqrt{\log(2+V)})$. This application uses the
finite Behrend construction inside the Lean proof and the density
Hales--Jewett theorem from outside it; neither estimate requires RH.

\paragraph{Comparison with a local algebra.}
For even $q$ let $A_b=\mathbb F_q\oplus V$ with $\dim_{\mathbb F_q}V=b-1$
and $V^2=0$. The earlier
\href{https://github.com/VladimirReshetnikov/ProveIt/blob/5ed40f562710fcc2770bbb463dba528a6727f7ae/docs/incoming/grh_method_transfers_research.zip}{ProveIt report},
in its section on finite local rings, proves
$1-Q^{-2}\le\Delta(A_b)\le1-(Q^2+Q+1)^{-1}$ for $Q=q^b$.
One can see the limit $1$ already from the earlier report's
squarefree-gcd construction: in this algebra $(a+v)^2=a^2$, so a
nonunit square producing a nonzero endpoint forces every scalar component
of that endpoint to have a common squared prime divisor. By excluding
all such endpoints one retains natural density $1-q/Q^2$.
For the same $q,b$, the reduced algebra $\mathbb F_q^b$ has the same
cardinality $Q$, yet its maximal upper density tends to zero. Its
coordinate idempotents permit the independent ratios that realize the
ternary word cube. Question~6.5 includes both examples; its general
finite-ring form and the exact finite-factor extrema remain open here.
